\chapter{Requirements Specifications}\label{chap:req}

The limitations discussed in Subsection~\ref{sub:gaps} have shaped the technical vision for the proposed system. To address these gaps and deliver a robust, 
scalable educational platform following our goal (Section~\ref{sec:goal}), this section outlines the explicit Functional and Non-Functional Requirements, as well as 
visiallizing use---case design. These requirements are derived directly from the systemic constraints and the pedagogical needs of learners for a more personalized, 
diversed and comprehensive learning assistance system.

\section{Functional Requirements}

\subsection{Knowledge Modelling and Management}
As a key objective feature, knowledge modelling construct knowledge base, foundation to all other key features. This functional category encompasses the capabilities 
required to process raw educational inputs into a structured knowledge base.

\begin{itemize}
    \item \textbf{Data Ingestion and Parsing:} 
    \begin{itemize}
        \item The system must be able to ingest unstructured, text-dense academic formats (e.g., PDF textbooks, syllabus guidelines).
        \item The system must accurately parse text into discrete chunks while retaining crucial positional metadata (e.g., source file origin, exact page numbers).
    \end{itemize}
    
    \item \textbf{Semantic Extraction and Normalization:}
    \begin{itemize}
        \item The system must extract fundamental educational concepts and their explicit pedagogical relationships from the ingested text.
        \item The system must normalize extracted entities to a canonical form to ensure conceptual consistency and prevent any form of semantic fragmentation (even of different course) 
        across the knowledge base.
    \end{itemize}
    
    \item \textbf{Hierarchical and Relational Organization:}
    \begin{itemize}
        \item The system must organize extracted entities into a standardized hierarchical topology.
        \item The system must automatically group highly related concepts into cohesive educational clusters or chapters.
        \item The system must explicitly define and store directional "prerequisite" relationships between concepts to enable automated learning path generation.
    \end{itemize}
    
    \item \textbf{Storage Capability:}
    \begin{itemize}
        \item The persistence mechanism must guarantee a database structure that facilitates high-speed, complex relational querying and subgraph retrieval during live 
        interactions.
        \item To ensure diversity of learning input, storage design must ensure efficient storing of enormous data with reasonable latency. 
    \end{itemize}
\end{itemize}

\subsection{Intelligent Teaching Assistance}
To utilized modeled knowledge base as well as guarantee satisfactory learning experience and result, the design of agentic system that serves as our Teaching Assistance 
is crucial. This category defines the required behaviors and reasoning capabilities of the automated teaching assistant.

\begin{itemize}
    \item \textbf{Intent Recognition and Context Orchestration:}
    \begin{itemize}
        \item The system must classify user intents (e.g., Retrieve information, Generate roadmap, Clarify concepts) and autonomously route the request to the appropriate internal workflow.
        \item The system must dynamically gather necessary context before responding, which includes retrieving the student's personal mastery history and the relevant subset of the educational knowledge base.
    \end{itemize}
    
    \item \textbf{Intelligent Reasoning and Supervision:}
    \begin{itemize}
        \item The system must support a Socratic teaching interaction; it must be capable of generating probing questions to assess student comprehension 
        rather than merely providing direct answers.
        \item The system must objectively evaluate the student's response to determine whether to authorize progression to the next prerequisite topic.
        \item The system must be able to autonomously detect when it lacks sufficient context to answer a question and execute a fallback strategy that 
        synthesizes information from the broader knowledge base.
    \end{itemize}
    
    \item \textbf{Response Verifiability and Fallback:}
    \begin{itemize}
        \item To guarantee grounded responses, the system must prioritize explicit document referencing, citing exact source pages for its generated answers.
        \item The system must execute a broader knowledge synthesis fallback only when direct contextual concepts are missing from the immediate document.
    \end{itemize}
\end{itemize}

\subsection{Student Interaction and System Navigation}
This category defines the functional requirements concerning the end-user's experience, accessibility, and direct interaction with the intelligent system.

\begin{itemize}
    \item \textbf{Identity and Access Security:}
    \begin{itemize}
        \item The system must provide secure authentication capabilities (Login/Logout) to uniquely identify users.
        \item The system must maintain isolated, tamper-proof session states for each user to securely record individual learning progress without cross-user data exposure.
    \end{itemize}
    
    \item \textbf{Interactive Dialogue and Visualization:}
    \begin{itemize}
        \item The system must provide a natural language chat interface that supports continuous, multi-turn conversational interactions and real-time streaming responses.
        \item The system must dynamically render visual learning roadmaps, empowering students to make informed, user-centric choices regarding their educational progression rather than following a static curriculum.
        \item The interface must consolidate essential learning tools, providing seamlessly integrated functional components such as conversational history logs, a synchronized document viewer, and progress summarization dashboards.
    \end{itemize}
    
    \item \textbf{Advanced Navigation and Content Discovery:}
    \begin{itemize}
        \item The system must implement an intuitive, multi-level navigation structure (e.g., breadcrumb trails) to help students orient themselves within complex knowledge hierarchies.
        \item The system must incorporate a powerful semantic search utility, allowing users to rapidly query and locate specific educational materials, graph concepts, or historical interactions.
    \end{itemize}
    
    \item \textbf{Contextual and Responsive Feedback:}
    \begin{itemize}
        \item The system must be capable of emitting explicit interface triggers alongside its textual responses to automatically synchronize the user's interface (e.g., auto-scrolling the document viewer to the exact referenced page).
        \item Notifications and intelligent pedagogical prompts must be delivered in a non-intrusive, accessible manner that captures user attention while preserving the fluidity of the learning experience.
    \end{itemize}
\end{itemize}


\section{Non-Functional Requirements}
To create a powerful, reliable, and user-friendly educational platform, the Smart Learning System cannot stop at being functional, but need focusing on efficiency 
and overall system performance. Thus, it must also adhere to a set of non-functional requirements that govern its 
performance, scalability, security, and overall user experience.
\begin{itemize}
    \item \textbf{Scalability and Modularity:}
    \begin{itemize}
        \item The backend architecture must be designed as stateless, relying exclusively on external persistence layers for state management to ensure readiness 
        for future distributed scaling.
        \item System design must be modular, with clear separation of components reponsibility and dependency to facilitate future enhancements, maintenance, and potential integration of 
        additional features or third-party services.
    \end{itemize}

    \item \textbf{Resource Optimization and Latency Minimization:}
    \begin{itemize}
        \item Operating under the constraint of localized execution, the system must optimize memory and computational overhead, loading resource-heavy 
        models into memory only when explicitly demanded by an active workflow.
        \item All resource must be efficiently released immediately after use to minimize the system's memory footprint and ensure responsiveness.
        \item System must be scheduled to be constantly running, minimizing downtime and waiting time for any task.
        \item All resource must be managable, specifically for developers and system administrators, but crucial for maintaining system stability and performance.
    \end{itemize}

    \item \textbf{Reliability and Determinism:}
    \begin{itemize}
        \item The system must mitigate the inherent unpredictability of generative AI by enforcing strict, programmatic data schemas for all internal system communications and decision-making logic.
        \item Complex analytical tasks, such as calculating knowledge gaps, must be executed via deterministic programmatic logic rather than relying on probabilistic generative approximations.
    \end{itemize}

    \item \textbf{Observability and Traceability:}
    \begin{itemize}
        \item The system must maintain comprehensive telemetry across all automated workflows.
        \item Database state changes, including knowledge graph updates and student progress logs, must be monitored and saved to enable historical analysis and rollback capabilities.        
        \item Every automated decision, context retrieval action, and resource consumption metric must be actively logged to provide a transparent audit trail for both system debugging and pedagogical evaluation.
    \end{itemize}
    \item \textbf{User Experience and Interface Requirements:}
    \begin{itemize}
        \item User interface is focus on the simplicity and intuitiveness, ensuring that users of varying technical proficiency can navigate the system effectively.
        \item User interface must be responsive and accessible, adhering to modern web standards to ensure compatibility across a wide range of devices and screen sizes (phones, tablets, desktops).
        \item Client side shouldn't be pulled by heavy backend response latency. It should maintain a fluid, real-time interaction experience, minimizing latency 
        in response generation and interface updates. Mininizing \texttt{"user experience latency"} regardless of system performance 
        is a strategic design move to preserve conversational engagement as well as learning interest.
    \end{itemize}

\end{itemize}


